\begin{table*}[t]
\fontsize{6.8pt}{7.5pt}\selectfont
\centering
\caption{\textbf{Annotation codebook definitions.} The table defines the labels used to translate learning-science constructs into observable turn-level signals. \revise{Passive, active and constructive are user-engagement depth labels; emotional expression is a parallel descriptive label.} These labels are behavioural annotations of conversation text; they are not measures of durable learning, learner identity or instructional efficacy.}
\label{tab:annotation_codebook_definitions}
\setlength{\tabcolsep}{3pt}
\renewcommand{\arraystretch}{1.02}
\begin{tabularx}{\textwidth}{
>{\raggedright\arraybackslash}p{0.13\textwidth}
>{\raggedright\arraybackslash}p{0.15\textwidth}
>{\raggedright\arraybackslash}p{0.32\textwidth}
>{\raggedright\arraybackslash}X}
\toprule
Family & Label & Definition & Boundary rule \\
\midrule
User engagement & Passive (\texttt{P}) & Reception, acknowledgement or minimal continuation without substantial transformation. & Includes short confirmations or requests to continue; excludes task-relevant follow-up or reasoning. \\
User engagement & Active (\texttt{A}) & Task-relevant action, follow-up, application or manipulation of supplied information or task material. & The user takes a task-relevant next step, such as applying a suggestion, reformulating an artifact or requesting the next step, but does not test, revise or extend an idea. \\
User engagement & Constructive (\texttt{C}) & Higher-depth cognitive work in which the user generates, revises or extends ideas, explanations or constraints. & Requires explicit user reasoning or transformation; not inferred from long text alone. \\
\revise{Emotional descriptor} & \revise{Emotional expression (\texttt{E})} & \revise{Explicit emotion in the user turn, such as frustration, gratitude, anxiety or excitement.} & \revise{Coded in parallel with user engagement; not part of the passive--active--constructive cognitive-depth hierarchy.} \\
\midrule
Assistant support & Scaffolding & Assistant support that structures, diagnoses, regulates or redirects the user's process while leaving cognitive or task responsibility with the user. & Distinguished from reference responses that simply provide an answer, draft or deliverable. \\
Support intent & \texttt{I1} metacognitive & Support aimed at planning, monitoring, reflection or regulation of the user's work. & Captures the purpose of support, not its surface form. \\
Support intent & \texttt{I2} cognitive & Support aimed at concepts, mechanisms, task knowledge or problem-solving content. & Most scaffolded support in this corpus is cognitive in purpose. \\
Support intent & \texttt{I3} affective & Support aimed at encouragement, reassurance or affective regulation. & Used only when affective support is substantively present, not for polite tone alone. \\
\midrule
Support form & \texttt{M1} feedback & Evaluative information about the user's prior attempt, diagnosis or understanding. & Requires a target in the user's preceding contribution. \\
Support form & \texttt{M2} hinting & A cue, prompt or partial path that helps the user proceed without fully resolving the task. & Leaves the next inferential or task step with the user. \\
Support form & \texttt{M3} instructing & Explicit procedural guidance, ordered steps or strategy for doing the task. & More directive than hinting; may still count as scaffolding when it structures process. \\
Support form & \texttt{M4} explaining & A rationale, mechanism or causal account that clarifies why something works or fails. & Requires explanatory content, not merely a restatement of an answer. \\
Support form & \texttt{M5} modelling & A reusable pattern, worked example or demonstration intended to be adapted. & Distinguished from simply giving the final requested product. \\
Support form & \texttt{M6} questioning & A question that advances diagnosis, reflection, planning or problem solving. & Excludes generic clarification questions that do not support the user's process. \\
\bottomrule
\end{tabularx}
\vspace{1mm}
\begin{minipage}{0.96\textwidth}
\footnotesize
\textit{Note.} Supplementary Table~\ref{tab:prompt_template_summary} summarizes the prompt-template families used to apply these definitions, and Supplementary Table~\ref{tab:annotation_implementation_controls} reports the corresponding model, parser, retry and boundary-rule controls.
\end{minipage}
\end{table*}
